\section{The short series on the identity cover}

% BEGIN S76G1
\paragraph{S76G1. Constructing the section on the given moduli space.}
Use the literal identity cover of (ICL.5) with its given connection,
and put $K=c.F$, $M=l.M$, $S=\mathcal O_E^n$, $r=p+q$ and
$\mathcal Q=(\mathrm{LevQ}(D_w,l))^n$. The tuple pinned at the
upper level has character one and distinguished residue zero.
The lower tuple $x_0\in D_w.\mathrm{La}^n$ remains a variable.
In particular the weight in the short sum is
$w(q)=\zeta_M^{\mathrm{pairL}(x_0,\mathrm{levRep}(q))}$.
For $\Lambda$ unimodular and trivial at $1$ and $M$, choose odd $N$
with $M\mid N$ and $NS\subseteq I_\Lambda:=\Lambda\cap S$.
Take $x_0$ satisfying the specified $\mathrm{AdmX}$ and
$P\in\mathrm{harmL}(c,\Lambda,a)$. There is a section
$g\in\mathrm{shortGen}(X_j,c,\nabla,1,0,a,n,d)$ for precisely
these data.

Indeed set $m=[I_\Lambda:NS]^r$, $z_\Lambda=[\Lambda:NS]^r$ and
$\beta=\beta_H(\Lambda)\in K^\times$. In the definition of
$\beta_H$ either a nonzero witness is chosen or the value is one.
The numerator in A76T2 has exactly $z_\Lambda$ isotropic torsion
tuples for each $q$, while its denominator is
$\beta\,m\,|\mathcal Q|\,N^{\deg P}$. A76T2 supplies its unique
natural value on each framed point and its weight-$(a,n)$ section.
The identity pin has $\mathrm{ev}(c_\delta,y_2)=1$ on every field
lift. Thus the descended value is precisely the value required by
$\mathrm{IsShortRule}$, on every lifted pair in its definition.
The lattice, integer, tuple and polynomial just chosen give the
four witnesses for membership in $\mathrm{shortGen}$. Only one
lattice and polynomial enter this section.
% END S76G1

% BEGIN S76G2
\paragraph{S76G2. The finite Fourier sum and the lattice factor.}
Fix a good embedding $\sigma:K\hookrightarrow\mathbb C$, a reference
Hermitian space $X$, and its native level datum $D$ tied to the
identity cover. Fix $r\leq n$ with $n$ even, and take $\Lambda\in V_{n}(M)$, the original family of
self-dual lattices agreeing with $S$ at the level primes. Write
$\mathfrak m=D.\mathfrak m=\mathrm{Mfr}(l)$ and $L'=D.\mathrm{sh}(L)$.
Take the prescribed unit vacuum $v_0(L)$ for $S\otimes L'$,
where $L$ runs over the level family. For a fixed such $L$, take a
unit vacuum $v$ for $S\otimes L$, the marking and its dual, and
one $x_0$ adapted to this dual. Write $H(c)$ for the original
Heisenberg translation, and $p_A$ for the average with the vacuum
semicharacter of $A$. The tensor notation below means the source's
additive matrix lattice. Define the finite set
\[
 B=\big(\Lambda\otimes 2\overline{\mathfrak m}L\big)/
     \big(\Lambda\otimes 2\overline{\mathfrak m}\mathfrak m L\big).
\]
In the test $d(c)=\mathrm{datF}(D,\mathrm{Mfr}(l),c)$ the first
$r$ rows reduce to an allowed matrix $j\in D.J$ and contribute
$D.\lambda(j)$. Every remaining row is congruent to zero modulo
$\mathrm{Mfr}(l)$, with zero coefficient when either test fails.
Because $D.\mathfrak m=\mathrm{Mfr}(l)$ on this cover, these are
exactly the tests defining $D.\mathrm{dat}$. The two functions
$\phi_F$ and $D.\phi$ thus have the same summation quotient, each
coefficient and translate, and the same normalized projected vacuum
$\mathrm{unit}(p_{\Lambda\otimes L'}v_0(L))$.

For any allowed roots $Y(q)$ of the half-level points put
\[
 e(c,q)=\exp\!\left(-2\pi i\sum_b E_F(c_b,Y(q)_b)\right).
\]
The marked
half-level pairing and the adaptation of the first $r$ rows give
\[
   \sigma w(q)=t_1\sum_{c\in B}d(c)e(c,q),\qquad |t_1|=1 .
\]
Indeed on the identity cover $J_d=\mathcal O_E$. The level group
$D.J$ reduces to the identity modulo $\mathfrak m$. Hence the
nonzero coefficient occurs in the residue class with prescribed
first $r$ rows and zero remaining rows. At primes of $\mathfrak m$
the tensors reduce in the given row coordinates since $2$,
$\overline{\mathfrak m}$ and the denominator comparing $\Lambda$
with $S$ are units there. Elsewhere the quotient has no part.
The dual marking identifies the exponential on the indicated class
with the pinned power of $\zeta_M$; changing a representative adds
an integer to its exponent. The level tie allows the common phase
$t_1$ to be chosen before $\Lambda,N,Y$.

Put $p=p_{\Lambda\otimes L}v$. Finite character orthogonality gives
\[
 \sum_{q\in\mathcal Q}e(c,q) H(Y(q))p
   =|\mathcal Q|\,H(c)p_{\Lambda\otimes L'}p .
\]
Here the root translations enumerate cosets of the shifted lattice
modulo its intersection with the old one. On their common
sublattice $p$ already has the prescribed semicharacter.
Commuting $H(c)$ through the translates gives precisely $e(c,q)$.
Averaging over cosets therefore gives the shifted projector, and
using the sum in place of its average accounts for $|\mathcal Q|$.

For the remaining projection put
$s_\Lambda=[S:I_\Lambda]^{-r/2}>0$. One has
\[
 p_{\Lambda\otimes L'}p
     =t_2\,s_\Lambda\,\mathrm{unit}
                      (p_{\Lambda\otimes L'}v_0(L)),
\]
where $t_2$ depends on the two vacua at $L$ and has positive
modulus independent of $\Lambda,N,Y$. To see the separation,
the modification $S\rightsquigarrow\Lambda$ is trivial at the
level primes and $L\rightsquigarrow L'$ is trivial away from them.
In the finite projector averages the two changes can consequently
be taken in either order, by the Chinese remainder theorem on the
finite quotients. Changing the target lattice first projects $v$
onto the vacuum line spanned by $v_0(L)$, with coefficient $t_2$.
The squared overlap of two unit lattice vacua is the reciprocal
relative index: in the character average over one lattice, a
translate outside the other has inner product zero, by its
nontrivial commutator with an element of the other self-dual
lattice. The translates on the intersection contribute one.
This proves $|t_2|>0$. The norm of the subsequent source projection
is $s_\Lambda$, since the relative index of the tensor lattices
is $[S:I_\Lambda]^r$ by local freeness of the rank-$r$ factor.
Expanding the coefficient Fourier sum, exchanging the two finite
sums, and using the two projector identities now gives
\[
 \sum_q \sigma w(q)H(Y(q))p
     =t\,s_\Lambda\,D.\phi(L,\Lambda),
       \qquad t=t_1|\mathcal Q|t_2 .
\]
The factor $t$ is common to all allowed $\Lambda,N,Y$ at the
fixed marking and with the chosen vacua.
% END S76G2

% BEGIN S76G3
\paragraph{S76G3. Computing the reading at a unit frame.}
Use a supplied principal uniformization $U'$ at $\sigma$ with its
common marked period data and letter dictionary $e_L$. Fix one
component $k$. Use its supplied unitary mover $\gamma_k$ and
transport its marking to $L_k=\gamma_k\Lambda_k$. The given dual
marking and lower tuple satisfy the adaptation of S76G2 on $L_k$.
For a unit frame $x$ the uniformization reads the rigidified
theta generator as a holomorphic theta function $\theta$ of
weighted mean square one on the period torus. Transport to the
reference space preserves its periods, marking and weighted mean.
Write $x'=\gamma_k\operatorname{cast}(x)$, where the cast identifies
the embedding indices on the two spaces, and put
$P^\sigma=\mathrm{rename}_{e_L}(\mathrm{map}_{\sigma}P)$.

In the Fock trivialization, $\prod_{b=1}^n\theta(z_b)=
\kappa_0\,\mathrm{thA}(v,x',z)$ for the unit vacuum $v$ of
$S\otimes L_k$. Both sides have the period multipliers of the
same level-one theta characteristic in each variable, whose
space of sections is a line. The proportionality has
$|\kappa_0|=\rho_0>0$ independent of the unit frame on this
marked lattice. In symplectic period coordinates, write
$h=P_1^{-1}$ for the inverse first period block and
$\tau=hP_2$. The Gaussian mean-square integral of a theta
characteristic over a period parallelepiped is evaluated by
orthogonality in the real-period coordinates and then an ordinary
Gaussian integral in the others. The factors involving
$\det\operatorname{Im}\tau$ and $|\det h|$ combine as the
determinant of the frame Hermitian form. At a unit frame its
absolute Gram determinant in each embedding is one. The remaining
factors are the constants of the fixed polarization, vacuum and
symplectic coordinates. Positivity follows from the positive mean
and the nonzero normalized vacuum.

Evaluate the algebraic rule of S76G1 at this frame. Lift its $z_\Lambda$
isotropic torsion tuples by the even representatives used in the
finite vacuum sum. Their Heisenberg translates sum to
$z_\Lambda\,p_{\Lambda\otimes L_k}v$: the semicharacter discrepancies
are trivial on the even representatives and this is the projector
average multiplied by its enumeration cardinality. The marked
division points then act by $H(Y(q))$. The translated line sections
are read at $Nz$ in the analytic variable $z$.
Applying the harmonic jet in the letter directions removes the
quadratic change of frame as in S76G1. The chain rule contributes
$N^{\deg P}$ because the monomials have the same total degree.
This cancels precisely the power in the denominator of the rule.
Linearity in the finite test and in $P$, followed by S76G2, gives
\[
 \mathrm{rd}_{U'}(g)_k(x)
    =\kappa_x u_{\Lambda,N}\,\mathrm{thP}
          (D.\phi(L_k,\Lambda),P^\sigma,x'),\qquad
 u_{\Lambda,N}=
    \frac{[\Lambda:NS]^r}{\sigma(\beta_H(\Lambda))
                    [I_\Lambda:NS]^r|\mathcal Q|}\,s_\Lambda .
\]
Here $\kappa_x=\kappa_0t$ since the identity-cover section
$c_\delta$ evaluates to one. The finite cardinality in $t$ is
the one coming from the Fourier sum of S76G2.
At a given unit frame the same $\kappa_x$ works simultaneously
for all admissible lattices, integers, polynomials and sections
obeying the rule. Its modulus is fixed and positive on this
component. The lattice factor is nonzero and can be chosen before
the component: the displayed index ratio is $[\Lambda:I_\Lambda]^r$
for every admissible $N$.
% END S76G3

% BEGIN S76G4
\paragraph{S76G4. One comparison scalar on a whole component.}
Index by every allowed quadruple $(\Lambda,N,P,g)$, with $g$ of
weight $(a,n)$ satisfying the original short rule. In the formula
of S76G3 write $R_i$ for its reading and $T_i$ for the entire
theta expression including $u_{\Lambda,N}$ but excluding $\kappa_x$.
Both are holomorphic on the component frame domain. Their Borel
character agrees by the weight and letter dictionary, including
the pullback by $\gamma_k$. In each embedding, triangular
Gram--Schmidt puts an arbitrary domain frame into unit form.
The common Borel character is nonzero, so cancellation transports
the simultaneous equality $R_i(x)=\kappa_x T_i(x)$ to the original
frame, still with prescribed positive modulus.

If every $T_i$ vanishes on the domain, every $R_i$ vanishes and any
scalar of that modulus works. Otherwise choose $T_j(x_*)\ne0$.
On a ball where $T_j$ is nonzero, $R_j/T_j$ is holomorphic with
constant modulus. Restriction to complex lines and the maximum
modulus principle give one constant $C_k\ne0$ there. The pointwise
equalities imply $R_i=C_kT_i$ on this ball for every $i$. The frame
domain is connected: per embedding the orthonormal frames are a
torsor for the connected real unitary group of the given signature,
and the invertible complex triangular factors form a connected
group. The identity theorem for each $R_i-C_kT_i$ extends the same
equality throughout the component. Thus (ICL.12) holds on these
charts for the actual single-lattice sections, with one $C_k$
for the entire family and the explicit common $u_{\Lambda,N}$.

In the spanning application use the system at the auxiliary
integer $2\,\mathrm{bad}\,M$ (here $\mathrm{bad}=1$). Its
lattice condition gives unimodularity, triviality at the two
required moduli and an allowed odd multiple $N$ of $M$ for each
$\Lambda$. The finite integral harmonic expansion of ICL15 supplies
the $P$ of S76G1. The same $x_0$ must be adapted on every component
and meet $\mathrm{AdmX}$ for these generators. For a simultaneous
expansion of the component readings divided by $C_k$, division
of each coefficient by $u_{\Lambda,N}$ now replaces its theta
term with the reading of $g$ on every component. The source
readings vanish off the frame domain, and ICL8 then gives coefficient
descent whenever the initial reading belongs to an original section.
This establishes the short-section comparison in ICL16 on supplied
marked charts with their tied level and normalized vacuum data.
% END S76G4
